\chapter{Degenerative Diseases}

\section*{Patient focus}
As we age, our joints and spine naturally undergo wear and tear. Imaging can show these changes---thinning of cartilage, wear on spinal discs, bone spurs, and narrowing of the spaces where nerves travel. However, many people with significant wear on their scans have few symptoms, while others with mild changes can have severe pain. The most important point is that imaging findings must be interpreted in the context of your symptoms. Not every abnormality on a scan needs treatment.


\label{ch:ch69_degenerative}

\section{Imaging Role}
Degenerative findings are common and may be incidental. Imaging should be ordered when the result is likely to change management or when red flags suggest an alternative diagnosis.

\begin{pearlbox}[Order imaging when it changes management]
Degenerative findings increase with age---disc bulges, osteophytes, and facet arthropathy are present in most asymptomatic adults over 50. Imaging is indicated when: (1) symptoms persist despite conservative treatment, (2) red flags are present, (3) surgical planning is needed, or (4) serial assessment of progression is required.
\end{pearlbox}

\section{Osteoarthritis Imaging}

\subsection{Kellgren-Lawrence Grading (Knee)}
The Kellgren-Lawrence (KL) system is the standard radiographic grading for knee OA:

\begin{table}[h]
\centering
\caption{Kellgren-Lawrence grading system for knee osteoarthritis}
\begin{tabular}{|l|l|l|}
\hline
\textbf{Grade} & \textbf{Description} & \textbf{Key features} \\
\hline
0 & Normal & No radiographic findings \\
1 & Doubtful & Minute osteophytes, questionable joint space narrowing \\
2 & Minimal & Definite osteophytes, possible joint space narrowing \\
3 & Moderate & Moderate osteophytes, definite joint space narrowing, some sclerosis \\
4 & Severe & Large osteophytes, marked joint space narrowing, severe sclerosis, deformity \\
\hline
\end{tabular}
\end{table}

\subsection{Radiographic Features of OA}
\begin{itemize}
  \item \textbf{Joint space narrowing}: asymmetric (medial compartment in knee; superior in hip)
  \item \textbf{Osteophytes}: bony projections at joint margins---the hallmark of OA
  \item \textbf{Subchondral sclerosis}: increased bone density beneath cartilage
  \item \textbf{Subchondral cysts}: well-defined lucent areas in subchondral bone
  \item \textbf{Deformity}: varus/valgus alignment in knee, femoral head migration in hip
\end{itemize}

\subsection{Modality Choice in OA}
\begin{table}[h]
\centering
\caption{Imaging in osteoarthritis}
\begin{tabular}{|l|l|l|}
\hline
\textbf{Modality} & \textbf{Role} & \textbf{Advantages} \\
\hline
CT & Complex fractures, pre-op 3D planning & Bony detail, alignment \\
MRI & Early OA, internal derangement, pre-op & Cartilage, meniscus, ligaments \\
Ultrasound & Effusion, Baker's cyst, injection guidance & Real-time, no radiation \\
X-ray & First-line, staging, surgical planning & Quick, cheap, KL grading \\
\hline
\end{tabular}
\end{table}

\section{Degenerative Disc Disease}

\subsection{MRI Findings}
\begin{itemize}
  \item \textbf{Normal disc}: high T2 signal (hydrated nucleus pulposus), intact annulus
  \item \textbf{Desiccation}: low T2 signal in nucleus---early degenerative change
  \item \textbf{Disc height loss}: progressive finding
  \item \textbf{Disc bulge}: symmetric extension of disc material beyond vertebral body margins
  \item \textbf{Disc protrusion (herniation)}: focal extension beyond disc space; may be contained
  \item \textbf{Disc extrusion}: material extends through annulus but remains connected to parent disc
  \item \textbf{Sequestration}: fragment completely detached from parent disc
\end{itemize}

\subsection{Pfirrmann Grading (Lumbar Disc)}
\begin{table}[h]
\centering
\caption{Pfirrmann grading of disc degeneration on MRI}
\begin{tabular}{|l|l|l|}
\hline
\textbf{Grade} & \textbf{T2 signal} & \textbf{Disc structure} \\
\hline
I & Homogeneous bright & Normal \\
II & Bright with horizontal band & Inhomogeneous, unclear disc/AF border \\
III & Intermediate grey & Inhomogeneous, unclear disc/AF border \\
IV & Grey to dark & Inhomogeneous, lost disc/AF distinction \\
V & Dark black & Complete disc space collapse \\
\hline
\end{tabular}
\end{table}

\section{Spinal Stenosis}

\subsection{Lumbar Spinal Stenosis}
\begin{itemize}
  \item \textbf{Central stenosis}: AP diameter $<$10 mm; caused by disc bulge, facet hypertrophy, ligamentum flavum thickening, and osteophytes
  \item \textbf{Foraminal stenosis}: narrowing of neural foramen by osteophytes or disc---causes radiculopathy
  \item \textbf{Lateral recess stenosis}: narrowing of the lateral recess before the foramen---compresses traversing nerve root
  \item \textbf{MRI is gold standard}: shows disc, facet, ligamentum flavum, and nerve root compression in one study
  \item \textbf{Sagittal grading} (Schizas): A (CSF visible around cauda equina) to D (no CSF visible---severe stenosis)
\end{itemize}

\subsection{Cervical Spinal Stenosis}
\begin{itemize}
  \item \textbf{AP canal diameter}: normal $>$13 mm, stenosis $<$10 mm
  \item \textbf{Causes}: disc osteophyte complex, ossification of posterior longitudinal ligament (OPLL), facet hypertrophy
  \item \textbf{Myelopathy signs on MRI}: T2 high signal in cord (oedema), cord compression, cord atrophy
  \item \textbf{T2 cord signal}: indicates myelopathy and may predict surgical outcome---earlier intervention is better
\end{itemize}

\begin{pearlbox}[Stenosis---symptoms don't always match scans]
Severe spinal stenosis on MRI may be asymptomatic. Conversely, mild stenosis in a narrow canal may cause significant claudication. Always correlate with clinical symptoms: neurogenic claudication (stopping with rest, relieved by flexion) suggests lumbar stenosis; myelopathic signs (upper motor neuron signs, gait disturbance) suggest cervical myelopathy.
\end{pearlbox}

\section{Spondylolisthesis}

\subsection{Types and Grading}
\begin{itemize}
  \item \textbf{Degenerative (Type III)}: most common in elderly; slip at intact neural arch---facet joint arthropathy
  \item \textbf{Isthmic (Type II)}: pars interarticularis defect; most common in young adults
  \item \textbf{Grading} (Meyerding): percentage of vertebral body slippage over the one below:
    \begin{enumerate}
      \item Grade I: $<$25\% slip
      \item Grade II: 25--50\% slip
      \item Grade III: 50--75\% slip
      \item Grade IV: 75--100\% slip
      \item Grade V (spondyloptosis): $>$100\% slip
    \end{enumerate}
\end{itemize}

\subsection{Imaging}
\begin{itemize}
  \item \textbf{X-ray (lateral)}: first-line; measure slip percentage, assess alignment
  \item \textbf{Flexion-extension views}: assess stability (dynamic slip $>$3 mm or $>$15\% suggests instability)
  \item \textbf{MRI}: assess nerve compression, disc degeneration at slip level, and adjacent segments
  \item \textbf{CT}: best for pars defect and bony detail
\end{itemize}

\section{Osteoporosis Imaging}

\subsection{DEXA (Dual-Energy X-ray Absorptiometry)}
\begin{itemize}
  \item \textbf{Gold standard} for bone mineral density (BMD) measurement
  \item \textbf{Sites}: lumbar spine (L1--L4) and femoral neck
  \item \textbf{T-score}:
    \begin{enumerate}
      \item $\geq$-1.0: normal
      \item -1.0 to -2.5: osteopenia
      \item $\leq$-2.5: osteoporosis
    \end{enumerate}
  \item \textbf{Fracture risk}: FRAX tool combines T-score with clinical risk factors
\end{itemize}

\subsection{Vertebral Fracture Assessment (VFA)}
\begin{itemize}
  \item Low-dose lateral spine imaging (performed on DEXA machine)
  \item Identifies vertebral compression fractures (morphometric and non-morphometric)
  \item Presence of vertebral fracture automatically classifies as osteoporosis regardless of T-score
  \item Grading: Genant classification (Grade 1: 20--25\% height loss; Grade 2: 25--40\%; Grade 3: $>$40\%)
\end{itemize}

\subsection{CT and MRI in Osteoporosis}
\begin{itemize}
  \item \textbf{CT}: detects fractures not visible on X-ray; Hounsfield unit measurement of vertebral bodies (HU $<$110 suggests osteoporosis)
  \item \textbf{MRI}: distinguishes acute (oedema, high T2/STIR signal) from chronic (fatty marrow, low signal) vertebral fractures---essential when malignancy needs exclusion
  \item \textbf{Bone scan}: increased uptake in fracture---but non-specific
\end{itemize}

\begin{pearlbox}[Acute vs chronic vertebral fracture]
Acute osteoporotic fracture: high T2/STIR signal (oedema), normal posterior elements, preserved height (or recent loss). Malignant fracture: destruction of posterior elements, soft tissue mass, diffuse marrow replacement, loss of posterior vertebral body height. When in doubt, MRI with contrast is the key differentiator.
\end{pearlbox}

\section{Interpretation}
Describe distribution, severity, alignment, joint or disc-space change, marrow reaction, stenosis, nerve or cord effect, and progression. Correlate with symptoms because structural severity and clinical disability may not match.

\section{Reporting}
Avoid attributing symptoms to common findings without clinical support. Highlight fracture, infection, tumor, inflammatory disease, or severe neural compromise when present, and give proportionate recommendations.

\begin{pearlbox}[Reporting degenerative disease]
A good report distinguishes:
\begin{enumerate}
  \item \textbf{Incidental degenerative findings}: present in most adults; note but do not overcall
  \item \textbf{Findings that may explain symptoms}: correlation with clinical presentation
  \item \textbf{Red flags}: fracture, infection, tumour, inflammatory disease, cord compression
  \item \textbf{Actionable findings}: high-grade stenosis with myelopathy, unstable spondylolisthesis, large extrusion with cauda equina syndrome
\end{enumerate}
Use proportionate language---``degenerative changes'' for mild findings; ``moderate central canal stenosis with cord compression and T2 signal'' for significant findings requiring urgent communication.
\end{pearlbox}

\clearpage
